\documentclass[11pt]{article}
\usepackage[utf8]{inputenc}
\usepackage{amsmath,amssymb,amsthm}
\usepackage[margin=1in]{geometry}
\usepackage{microtype}
\usepackage{hyperref}
\usepackage{xcolor}
\usepackage{enumitem}
\usepackage{newunicodechar}
\emergencystretch=3em
\newunicodechar{ℝ}{\ensuremath{\mathbb{R}}}
\newunicodechar{ℕ}{\ensuremath{\mathbb{N}}}
\newunicodechar{∫}{\ensuremath{\int}}
\newunicodechar{λ}{\ensuremath{\lambda}}
\newunicodechar{α}{\ensuremath{\alpha}}
\newunicodechar{γ}{\ensuremath{\gamma}}
\newunicodechar{∂}{\ensuremath{\partial}}
\newunicodechar{²}{\ensuremath{^2}}
\newunicodechar{³}{\ensuremath{^3}}
\newunicodechar{·}{\ensuremath{\cdot}}
\newunicodechar{≤}{\ensuremath{\leq}}
\newunicodechar{₀}{\ensuremath{_0}}

\newcommand{\deriv}{\operatorname{deriv}}
\newcommand{\tideref}[2][]{\href{https://therisensea.org/tides/#2/}{\ifx&#1&\texttt{#2}\else#1\fi}}

\title{Tide retrospective:\\\emph{general-\(h\) partition derivatives and the formal \(Z''(0), Z'''(0)\)}}
\author{Claude Opus 4.8 (1M ctx), with Daniel Murfet}
\date{2026-05-30}

\begin{document}
\maketitle

\begin{abstract}
The previous tide proved the iterated-derivative \emph{step} for the
anharmonic perturbed partition family $G_n(h)=\int (-(tx))^n e^{-t(L+hx)}$
only at $h=0$, which gives the derivative of each family member but does not
formally identify the higher derivatives of $Z=G_0$. Following GPT-5.5 Pro's
recommendation, this tide proves the \textbf{local-in-$h$} step --- for every
$n$ and $|h|<1$, $\frac{d}{dh}G_n(h)=G_{n+1}(h)$ --- and chains it through
\texttt{iteratedDeriv} to land the genuine identities
\[
  Z''(0)=\int (tx)^2 e^{-tL},\qquad Z'''(0)=\int -(tx)^3 e^{-tL}.
\]
New file \texttt{AnharmonicPartitionDerivGeneralH.lean}; commit \texttt{a7886cd}
on branch \texttt{tide/anharmonic-partition-deriv-general-h}; full library
builds (8304 jobs), no \texttt{sorry}/\texttt{axiom}/\texttt{native\_decide},
no lint warnings.
\end{abstract}

\section{Setting}

The laplace seabed formalises the Laplace/Gibbs calculations of the SLT
Susceptibility Primer. The prior tide\footnote{\tideref[the higher-derivative
step tide]{2026-05-30-tide-anharmonic-partition-higher-deriv}.} introduced the
weighted-partition family $G_n(h)=\int(-(tx))^n e^{-t(L+hx)}$ and proved
$G_n'(0)=G_{n+1}(0)$. In its post-merge consult, GPT-5.5 Pro flagged the gap:
proving the step only at $h=0$ shows that each $G_n$ has $G_{n+1}(0)$ as its
derivative at $0$, but to call $G_2(0)$ ``$Z''(0)$'' one must know that
$\deriv Z = G_1$ \emph{on a neighbourhood of $0$} so the second derivative
can be taken. The fix it recommended --- prove the local-in-$h$ step, then
chain --- is this tide.

\section{The thing formalised}

\begin{quote}\itshape
For $\lambda,\gamma>0$, $\alpha^2<3\lambda\gamma$, $t>0$, every $n\in\mathbb{N}$
and every $h$ with $|h|<1$: $\ \mathrm{HasDerivAt}\ G_n\ \big(G_{n+1}(h)\big)\ h$.
Consequently $\deriv(G_n)(h)=G_{n+1}(h)$ for $|h|<1$,
$\ \mathrm{iteratedDeriv}\ k\ G_0\ h = G_k(h)$, and in particular
$Z''(0)=\int (tx)^2 e^{-tL}$, $Z'''(0)=\int -(tx)^3 e^{-tL}$.
\end{quote}

The headline Lean signatures:

\begin{verbatim}
theorem weightedPartition_hasDerivAt (n : ℕ) ... {h₀ : ℝ} (hh₀ : |h₀| < 1) :
    HasDerivAt (weightedPartition lam alpha gamma t n)
      (weightedPartition lam alpha gamma t (n + 1) h₀) h₀

theorem iteratedDeriv_two_partition_zero ... :
    iteratedDeriv 2 (weightedPartition lam alpha gamma t 0) 0
      = ∫ x, (t * x) ^ 2 *
          Real.exp (-(t * anharmonicPotential lam alpha gamma x))
\end{verbatim}

\section{Proof strategy}

\textbf{The step.} The dominated-differentiation argument is the prior tide's,
with the centre moved from $0$ to an arbitrary $h_0$ with $|h_0|<1$. Two
things change. First, the neighbourhood: $\mathrm{ball}\,0\,1$ is open and
contains $h_0$, hence is a neighbourhood of $h_0$, so the same ball and the
same $|h|<1$ bound serve at any interior centre. Second, the
dominated-differentiation lemma needs integrability of the family
\emph{at the centre} $h_0$; with $h_0\neq0$ the $+0\cdot x$ simplification of
the $h=0$ case is gone, so this needs the bare-exponential bound
$e^{-t(L+h_0x)}\le e^{t/2c}e^{-(t/2)L}$ ($|h_0|<1$), which I extracted as
\texttt{anharmonic\_perturbed\_exp\_le} (the exponent half of the existing
weighted pointwise bound; \texttt{amgm\_t\_abs\_x} un-privated to support it).

\textbf{The chaining.} With $\deriv(G_n)(h)=G_{n+1}(h)$ on $\mathrm{ball}\,0\,1$
in hand, an induction on $k$ gives $\mathrm{iteratedDeriv}\,k\,G_0\,h=G_k(h)$
for $|h|<1$: at the successor, the inductive hypothesis holds on the whole
open ball, hence $\mathrm{iteratedDeriv}\,k\,G_0$ agrees with $G_k$ on a
neighbourhood of $h$, so \texttt{Filter.EventuallyEq.deriv\_eq} turns
$\mathrm{iteratedDeriv}(k{+}1)\,G_0\,h=\deriv(\mathrm{iteratedDeriv}\,k\,G_0)\,h$
into $\deriv(G_k)(h)=G_{k+1}(h)$. The $k=2,3$ specialisations at $h=0$,
with the integrand simplified by \texttt{ring} (treating the exponential as an
atom), are the headline $Z''(0)$, $Z'''(0)$.

\section{Roadblocks and resolutions}

The chaining went in cleanly once phrased as ``agree on the open ball'' rather
than ``agree at a point'': it is the open-ball membership that supplies the
\texttt{EventuallyEq} the derivative-rewrite needs, and stating the induction
hypothesis as $\forall h,\ |h|<1\to\dots$ rather than a single
\texttt{EventuallyEq} made the successor step's neighbourhood argument a
two-line \texttt{Filter.eventuallyEq\_of\_mem}.

The one build failure was cosmetic: the final integrand-matching used
\texttt{congr 1; ring} on $(-(tx))^k = \pm(tx)^k$, but \texttt{congr} tries to
peel the power and strands an impossible base equality. Rewriting the
exponential factor first and then closing with a single \texttt{ring} (the
exponential is an opaque atom, the rest is a polynomial identity) fixed both
the $k=2$ and $k=3$ cases.

\section{What was Mathlib, what was new}

\textbf{Mathlib.} The dominated-differentiation
lemma\footnote{\texttt{hasDerivAt\_integral\_of\_dominated\_loc\_of\_deriv\_le}.};
\texttt{HasDerivAt.deriv}; \texttt{iteratedDeriv\_zero},
\texttt{iteratedDeriv\_succ}; \texttt{Filter.EventuallyEq.deriv\_eq},
\texttt{Filter.eventuallyEq\_of\_mem}; \texttt{Metric.isOpen\_ball.mem\_nhds}.

\textbf{Seabed (reused).} The general $|x|^m$ integrability witness, the two
public pointwise helpers, and \texttt{anharmonic\_coercive} --- all from the
two earlier anharmonic tides this day; \texttt{amgm\_t\_abs\_x} (un-privated).

\textbf{New.} The bare-exp domination bound; integrability at a nonzero
centre; the general-$h$ step; the \texttt{iteratedDeriv} induction; and the
two genuine partition-derivative identities. About 250 lines.

\section{Lessons}

\begin{itemize}[itemsep=3pt]
\item \textbf{A post-merge consult is still worth acting on.} GPT's
  ``prove the local-in-$h$ version'' arrived after the $h=0$ tide had merged;
  treating it as the next tide rather than ignoring it turned a
  derivative-of-a-family-member into the genuinely useful $Z''(0)$, $Z'''(0)$.
\item \textbf{For \texttt{iteratedDeriv} chaining, carry the hypothesis on an
  open set, not a point.} The neighbourhood agreement that
  \texttt{EventuallyEq.deriv\_eq} consumes falls straight out of open-ball
  membership; phrasing the induction over $|h|<1$ avoids any bespoke
  \texttt{EventuallyEq} bookkeeping.
\item \textbf{\texttt{ring} with an opaque \texttt{exp} atom beats \texttt{congr}
  for integrand matching.} \texttt{congr} peels powers and strands false base
  goals; rewriting the exponential then \texttt{ring} closes signed-power
  identities in one step.
\end{itemize}

\section{Follow-ups}

\begin{itemize}[itemsep=3pt]
\item \textbf{Fourth-cumulant / flow-equation identity.} With $Z(0)>0$ and
  $Z',Z'',Z''',Z^{(4)}$ at $0$ all available (the last is the $k=4$
  specialisation), assemble the connected fourth cumulant
  $\kappa_4=\partial_h^4\log Z|_0$. This is the motivating application and is
  now fully unblocked.
\item \textbf{Moment normalisation.} $Z^{(n)}(0)/Z(0)=\langle(-tx)^n\rangle_0$
  packages the derivatives as raw moments under the unperturbed Gibbs measure.
\item \textbf{Promote \texttt{amgm\_t\_abs\_x} to \texttt{lean-common}.} Now
  load-bearing across three anharmonic tides and genuinely generic.
\end{itemize}

\end{document}
